\section*{Chapter \ref{ch:g10:exercise-and-energy} --- Exercise and the Body's Energy Needs}

\begin{solution}{exo:g10:exercise-and-energy:1}
The energy spent at complete rest to keep the body alive --- heart,
brain, temperature, \omterm{def:g10:cells-common-unit:cell}{cell} maintenance: about \qty{80}{W}, some
\qty{7000}{kJ} a day.
\end{solution}

\begin{solution}{exo:g10:exercise-and-energy:2}
$2.0 \times 20 = \qty{40}{kJ}$ per minute, i.e. $40000/60 \approx
\qty{670}{W}$.
\end{solution}

\begin{solution}{exo:g10:exercise-and-energy:3}
The muscle's phosphate store (about ten seconds); lactic \omterm{def:g10:cell-metabolism:fermentation}{fermentation}
(one to two minutes of all-out effort); \omterm{prop:g10:cell-metabolism:respiration}{cellular respiration} (hours).
\end{solution}

\begin{solution}{exo:g10:exercise-and-energy:4}
The highest rate at which the body can take up and use oxygen.
$3500/70 = \qty{50}{mL/min}$ per kilogram.
\end{solution}

\begin{solution}{exo:g10:exercise-and-energy:5}
Muscles, about \qty{400}{g}, fuelling the muscle that stores it;
liver, about \qty{100}{g}, released as glucose into the blood for the
brain and the other tissues.
\end{solution}

\begin{solution}{exo:g10:exercise-and-energy:6}
\qty{2.7}{L/min}: $2.7 \times 20 = \qty{54}{kJ}$ per minute, i.e.
\qty{900}{W}. Heat: $900 - 200 = \qty{700}{W}$.
\end{solution}

\begin{solution}{exo:g10:exercise-and-energy:7}
Work $mgh = 60 \times 9.8 \times 800 \approx \qty{470}{kJ}$. At 25\%
efficiency the muscles spent about \qty{1900}{kJ}. Over \qty{7200}{s}:
about \qty{260}{W} above rest.
\end{solution}

\begin{solution}{exo:g10:exercise-and-energy:8}
A 100-metre sprint runs almost entirely on the phosphate store: no
oxygen debt, little lactic acid. A 400-metre race lasts about a
minute and draws half its energy from lactic \omterm{def:g10:cell-metabolism:fermentation}{fermentation}: lactic acid
accumulates (the ache) and the respiration that must clear it keeps
the runner breathing hard for minutes afterwards.
\end{solution}

\begin{solution}{exo:g10:exercise-and-energy:9}
\omterm{def:g10:chemistry-of-life:families}{Carbohydrate} energy $0.6 \times 3000 = \qty{1800}{kJ}$; \omterm{def:g10:exercise-and-energy:glycogen}{glycogen}
$1800/17 \approx \qty{106}{g}$; water released about \qty{320}{g}.
\end{solution}

\begin{solution}{exo:g10:exercise-and-energy:10}
Metabolic power about \qty{600}{W}: $\num{1.1e6}/600 \approx
\qty{1830}{s}$, about 30 minutes.
\end{solution}

\begin{solution}{exo:g10:exercise-and-energy:11}
Three quarters of the energy released in the muscles is heat, several
hundred watts during hard exercise. Sweating removes it: evaporating
water takes about \qty{2.4}{kJ} per gram, and blood flow to the skin
carries the heat there. Without sweating, body temperature would rise
by a degree every few minutes.
\end{solution}

\begin{solution}{exo:g10:exercise-and-energy:12}
Uphill, the power needed is proportional to body mass: the
\qty{60}{kg} subject has \qty{60}{mL/min} per kilogram against
\num{40}, and climbs faster. On the laboratory bicycle the power
does not depend on mass, both have the same \qty{3.6}{L/min}, and
both can sustain the same watts.
\end{solution}

\begin{solution}{exo:g10:exercise-and-energy:13}
The missing \qty{0.4}{L/min} of oxygen-equivalent is supplied by
lactic \omterm{def:g10:cell-metabolism:fermentation}{fermentation} in the muscles. Lactic acid accumulates; within a
few minutes the muscles burn and the runner is forced to slow to a
pace that respiration alone can cover.
\end{solution}

\begin{solution}{exo:g10:exercise-and-energy:14}
Short maximal bursts run on the phosphate store and \omterm{def:g10:cell-metabolism:fermentation}{fermentation},
which use no fat at all; fat is burned only by respiration, which
takes minutes to reach full rate and dominates only in prolonged
moderate exercise. The total energy of a few bursts is also small
compared with an hour's steady effort.
\end{solution}

\begin{solution}{exo:g10:exercise-and-energy:15}
Deficit $25000 - 8500 - 6000 = \qty{10500}{kJ}$, i.e. about
\qty{280}{g} of fat. Eating during the stage keeps the blood glucose
up for the brain and spares the \omterm{def:g10:exercise-and-energy:glycogen}{glycogen} for the climbs, where the
power needed is far above what fat alone can supply; the fat covers
the steady part of the day.
\end{solution}

\begin{solution}{pb:g10:exercise-and-energy:1}
\textbf{1.} $4.0 \times 60 \times 42.2 \approx \qty{10100}{kJ}$.

\textbf{2.} $\qty{3}{h}\,\qty{30}{min} = \qty{12600}{s}$:
$\num{1.01e7}/12600 \approx \qty{800}{W}$.

\textbf{3.} $10100/20 \approx \qty{506}{L}$; over \qty{210}{min}, about
\qty{2.4}{L/min}.

\textbf{4.} $2400/60 = \qty{40}{mL/min}$ per kilogram, i.e. $40/52
\approx 77\%$ of her maximum.

\textbf{5.} $(2.4 - 0.2)/2.4 \approx 92\%$.

\textbf{6.} $380 \times 17 = \qty{6460}{kJ}$.

\textbf{7.} The race costs \qty{240}{kJ} per kilometre: $6460/240
\approx \qty{27}{km}$.

\textbf{8.} \omterm{def:g10:chemistry-of-life:families}{Carbohydrate} use $0.8 \times 240 = \qty{192}{kJ/km}$:
$6460/192 \approx \qty{34}{km}$ --- the wall.

\textbf{9.} Fat energy $0.2 \times 10100 \approx \qty{2020}{kJ}$, i.e.
$2020/38 \approx \qty{53}{g}$: about 0.6\% of her \qty{9}{kg}.

\textbf{10.} Fat is burned only by respiration and at a limited rate:
it can supply roughly half the power her pace requires. Without
\omterm{def:g10:exercise-and-energy:glycogen}{glycogen} the muscles cannot make ATP fast enough, and the pace falls.

\textbf{11.} $60 \times 3.5 = \qty{210}{g}$, i.e. $210 \times 17 \approx
\qty{3570}{kJ}$, covering $3570/192 \approx \qty{19}{km}$ of
\omterm{def:g10:chemistry-of-life:families}{carbohydrate} use.

\textbf{12.} \omterm{def:g10:chemistry-of-life:families}{Carbohydrate} available $6460 + 3570 = \qty{10030}{kJ}$:
$10030/192 \approx \qty{52}{km}$, beyond the finish. Yes.

\textbf{13.} $+\qty{200}{g}$ of \omterm{def:g10:exercise-and-energy:glycogen}{glycogen} $+ \qty{600}{g}$ of water:
\qty{0.8}{kg}, costing $0.8 \times 4.0 \approx \qty{3}{kJ}$ per
kilometre extra --- negligible against 240.

\textbf{14.} \omterm{def:g10:exercise-and-energy:glycogen}{Glycogen} $580 \times 17 = \qty{9860}{kJ}$ plus the drink:
about \qty{13400}{kJ}; \omterm{def:g10:chemistry-of-life:families}{carbohydrate} use $0.55 \times 240 =
\qty{132}{kJ/km}$: over \qty{100}{km}. The wall has moved far beyond
the finish, with a wide margin.

\textbf{15.} The brain burns only glucose, supplied by the liver's
\omterm{def:g10:exercise-and-energy:glycogen}{glycogen}. When that store is empty the blood glucose falls, and the
brain signals the emergency as dizziness and overwhelming fatigue,
whatever the muscles still hold.

\textbf{16.} $3.6 \times 55 \times 42.2 \approx \qty{8360}{kJ}$;
$8360/20 = \qty{418}{L}$ over \qty{125}{min}: about \qty{3.3}{L/min}.

\textbf{17.} $3340/55 \approx \qty{61}{mL/min}$ per kilogram, i.e.
$61/80 \approx 76\%$ --- the same fraction as Nadia. The ceiling
differs, not the fraction of it used.

\textbf{18.} Economy is the cost per kilogram per kilometre. For a
\qty{55}{kg} runner, $(4.0 - 3.6) \times 55 \times 42.2 \approx
\qty{930}{kJ}$ saved over the race, about 9\%.

\textbf{19.} \omterm{def:g10:exercise-and-energy:glycogen}{Glycogen} \qty{8500}{kJ}; \omterm{def:g10:chemistry-of-life:families}{carbohydrate} needed $0.7 \times
8360 \approx \qty{5850}{kJ}$: yes, the stores suffice without
drinking.

\textbf{20.} Without preparation Nadia's \omterm{def:g10:exercise-and-energy:glycogen}{glycogen} runs out around
kilometre 34; \omterm{def:g10:chemistry-of-life:families}{carbohydrate} loading before the race, drinking
\omterm{def:g10:chemistry-of-life:families}{carbohydrate} during it, and training that raises the share of fat
burned each push the wall beyond the finish line.
\end{solution}
